% SPDX-License-Identifier: GPL-3.0-or-later
% Copyright (C) 2026 Alessandro Greco
\chapter{Installation and torrc}\label{cap:17-torrc}
\aggiornatoal{23 September 2026}{tor 0.4.9.13; the \texttt{tor(1)} manual for the same version}

\begin{obiettivi}
By the end of this chapter, you will be able to:
\begin{itemize}[nosep]
  \item choose the appropriate way to install Tor for your purpose;
  \item read and write a \file{torrc} file and know where to find documentation for every option;
  \item configure tor as a client, with SOCKS ports, stream isolation, and a protected control port;
  \item validate a configuration before using it and diagnose problems from logs.
\end{itemize}
\end{obiettivi}

This chapter begins the practical part of the book. Previous chapters explained how the protocol works; now we learn to use and configure the program that implements it. Every option mentioned here is documented in the tor manual (\cmd{man tor}), which is the authoritative reference: defaults can change between versions.

% ---------------------------------------------------------------------------
\section{Which Tor Should You Install?}\label{sec:torrc-installare}

\begin{description}
  \item[Tor Browser] For web browsing, this is almost always the appropriate choice. It includes a preconfigured tor and a browser modified to avoid disclosing information that would undermine the network's protections (\cref{cap:23-browser}). Using tor with an ordinary browser, by contrast, leaves users open to many forms of identification.
  \item[The \texttt{tor} package] For using Tor with other applications or operating a relay or onion service. Linux distributions include it, although often not the latest version. The Tor Project maintains its own package repository for Debian and Ubuntu and distributes the \emph{Tor Expert Bundle}, containing tor and pluggable transports, for Windows and macOS.
  \item[From source] As in the book's laboratory (\cref{sec:rp-avvio}), for studying tor or obtaining an exact version.
  \item[Arti] The Rust client (\cref{cap:30-arti}), increasingly mature but configured differently.
\end{description}

Whichever method you choose, download the software only from the official site and verify its signature. A version of tor modified by an adversary would nullify its protections.

% ---------------------------------------------------------------------------
\section{The torrc File}\label{sec:torrc-formato}

tor reads its configuration from a text file, usually \file{/etc/tor/torrc} when it runs as a system service. The format is simple: one option per line, followed by its values. Lines beginning with \texttt{\#} are comments. Some options may appear more than once, such as multiple SOCKS ports. Options can also be supplied on the command line with two leading hyphens:
\begin{lstlisting}[language=bash,numbers=none]
$ tor --SocksPort 9150 --DataDirectory /tmp/tor-test
\end{lstlisting}
The manual documents hundreds of options; \cmd{tor --list-torrc-options} prints them all. Distributions may specify a second file of defaults using \texttt{--defaults-torrc}. The file specified with \texttt{-f} takes precedence.

\subsection{Check Before Use}

Three commands can prevent many surprises:
\begin{lstlisting}[language=bash,numbers=none]
$ tor --verify-config -f /etc/tor/torrc     # is the configuration valid?
Configuration was valid
$ tor --dump-config short -f /etc/tor/torrc # what differs from the defaults?
$ tor --list-modules                        # which modules are compiled?
relay: yes
dirauth: yes
dircache: yes
pow: yes
\end{lstlisting}
As we saw in \cref{cap:15-dos}, the last command reveals, for instance, whether onion service proof of work is available.

% ---------------------------------------------------------------------------
\section{A Client Configuration}\label{sec:torrc-client}

A client needs very little configuration. Indeed, the minimal configuration file is empty: by default, tor opens a SOCKS port at 9050 on localhost. Here is a more explicit, annotated example for a system tor service whose account can access \file{/var/lib/tor} and \file{/var/log/tor} (not a generic unprivileged invocation):
\begin{lstlisting}[language=torrc]
# Where tor stores state, keys, and cache
DataDirectory /var/lib/tor

# SOCKS port for applications (localhost only!)
SocksPort 127.0.0.1:9050

# Control port, protected by the cookie (see Chapter 18)
ControlPort 127.0.0.1:9051
CookieAuthentication 1

# Log at notice level
Log notice file /var/log/tor/notice.log
\end{lstlisting}

\begin{figure}[htbp]
\centering
% SPDX-License-Identifier: GPL-3.0-or-later
% Copyright (C) 2026 Alessandro Greco
%% Accessibility alternative text, used only in EPUB.
\altfig{Example tor architecture: browser and curl connect to SocksPort 9050 on loopback, while Stem and Nyx use ControlPort 9051 only if it is enabled and protected by cookie or password authentication. The local SOCKS and control protocols do not themselves encrypt traffic and their listeners should not be exposed to the network. DataDirectory is a filesystem directory containing keys, state, and cache, not a port. A relay, including a bridge, may expose ORPort 443 to its peers; the port number is illustrative, not mandatory.}
% Example ports and data directory of a tor process
\begin{tikzpicture}[font=\sffamily\scriptsize,
  porta/.style={draw,rounded corners=1pt,fill=torviola!10,minimum height=0.55cm,minimum width=2.9cm,font=\ttfamily\scriptsize},
  dir/.style={draw=torgrigio,rounded corners=1pt,fill=torgrigio!10,minimum height=0.55cm,minimum width=2.9cm,font=\ttfamily\scriptsize},
  app/.style={draw=torgrigio,fill=torgrigio!8,rounded corners=2pt,minimum height=0.6cm,minimum width=2.1cm}]
  \node[draw=torviola,fill=torviola!5,rounded corners=4pt,minimum width=3.6cm,minimum height=4.3cm] (tor) at (4.6,0) {};
  \node[font=\sffamily\small,text=torviola] at (4.6,1.85) {\textbf{tor process}};
  \node[porta] (socks) at (4.6,1.0) {SocksPort 9050};
  \node[porta] (ctrl) at (4.6,0.2) {ControlPort 9051};
  \node[porta] (or) at (4.6,-0.6) {ORPort 443};
  \node[dir] (dd) at (4.6,-1.5) {DataDirectory};
  \node[app] (br) at (0.6,1.0) {browser, curl};
  \node[app] (st) at (0.6,0.2) {Stem, Nyx};
  \node[app,fill=torverde!12,draw=torverde!60!black] (rete) at (9.0,-0.6) {Tor network};
  \draw[->] (br) -- (socks); \draw[<->] (st) -- (ctrl); \draw[<->] (or) -- (rete);
  \node[anchor=west,font=\sffamily\tiny,align=left,text=torgrigio] at (6.6,1.0) {localhost only:\\SOCKS not encrypted};
  \node[anchor=west,font=\sffamily\tiny,align=left,text=torgrigio] at (6.6,0.2) {if configured:\\cookie or password};
  \node[anchor=west,font=\sffamily\tiny,align=left,text=torgrigio] at (6.6,-1.5) {keys, guard state,\\consensus cache};
  \node[font=\sffamily\tiny,text=torgrigio,align=center] at (4.6,-2.5) {ORPort 443 is an example\\for relays and bridges only};
\end{tikzpicture}

\caption{Example ports and directory of a tor process. SOCKS on 9050 serves local applications; ControlPort 9051 is optional and must be authenticated. Neither listener should be exposed to the network. DataDirectory is a filesystem directory, not a port. ORPort 443 is one possible choice for a relay or bridge, not a required port number.}
\label{fig:porte}
\end{figure}

\begin{attenzione}[Do Not Expose the SOCKS Port]
The manual is explicit: the SOCKS protocol is neither encrypted nor authenticated. A \opt{SocksPort} listening on a network-accessible address turns the computer into an open proxy and exposes applications' requests in plaintext on the local network.
\end{attenzione}

\subsection{Stream Isolation}\label{sec:torrc-isolamento}

Multiple streams can share the same circuit. This is efficient, but it allows an exit to link different activities performed by the same user. If personal email and a pseudonymous forum share a circuit, the exit can tell that they come from the same person. tor therefore provides \textbf{stream isolation}: rules determining which streams may share a circuit. The rules are expressed as \opt{SocksPort} options:
\begin{description}
  \item[\texttt{IsolateClientAddr}] Streams from different client addresses cannot share circuits (enabled by default).
  \item[\texttt{IsolateSOCKSAuth}] Streams with different SOCKS credentials cannot share circuits (enabled by default).
  \item[\texttt{IsolateDestAddr}, \texttt{IsolateDestPort}] Isolation by destination.
  \item[\texttt{IsolateClientProtocol}] Isolation by protocol (SOCKS4, SOCKS5, DNS\ldots).
\end{description}
Different SOCKS ports are also isolated from each other by default: each listener is assigned a separate \emph{session group}. The \opt{SessionGroup} option can override this behavior and permit ports to share circuits. The most widely used mechanism is SOCKS authentication isolation: Tor Browser sends different SOCKS credentials for each visited website, obtaining a separate circuit for each site. We will observe this in \cref{cap:18-stem}.

A circuit is not used forever. Once \opt{MaxCircuitDirtiness} has elapsed---10 minutes by default, counted from first use---the circuit accepts no new streams, and tor builds another.

\subsection{Choosing Relays: Exercise Caution}

The \opt{EntryNodes}, \opt{ExitNodes}, and \opt{ExcludeNodes} options can force or exclude relays by fingerprint, nickname, address, or country code. \opt{StrictNodes} makes \opt{ExcludeNodes} mandatory even when doing so breaks functionality. These tools are useful for experiments and specific circumstances. Using them merely to \enquote{feel safer}, however, usually makes matters worse: narrowing relay selection reduces path diversity, makes the user stand out, and defeats the protections provided by weighted path selection and guards (\cref{cap:13-percorso}).

\subsection{Other Ports}

\begin{description}
  \item[\opt{DNSPort}] tor answers DNS requests, resolving names through the network (\cref{sec:stream-dns}).
  \item[\opt{TransPort}] Accepts \emph{transparent} TCP connections redirected by the system firewall. Systems such as Whonix use this mechanism to route application traffic through Tor (\cref{cap:23-browser}).
  \item[\opt{AutomapHostsOnResolve}] Assigns virtual addresses to \texttt{.onion} names requested through DNS, allowing applications that do not support SOCKS to reach onion services through \opt{TransPort}.
\end{description}

% ---------------------------------------------------------------------------
\section{The Control Port}\label{sec:torrc-controllo}

\opt{ControlPort} allows other programs to control tor while it is running: read its status, request new circuits, or create onion services dynamically. It is extremely powerful and must therefore be protected. Two authentication methods are available:
\begin{description}
  \item[Cookie] With \opt{CookieAuthentication} 1, tor writes a file containing 32 random bytes in its data directory. Anyone who can read that file can authenticate. Security therefore depends on file permissions.
  \item[Password] With \opt{HashedControlPassword}, tor stores a password hash generated by \cmd{tor --hash-password}:
\begin{lstlisting}[language=bash,numbers=none]
$ tor --hash-password test
16:458528FE0C0AF36260B5071446823FB789212185D7EE7F615BC120D55D
\end{lstlisting}
The hash shown is an illustrative, verified hash for the sample password \texttt{test}, using a fixed salt for the printed example. A real invocation chooses a random salt, so its output normally differs even for the same password. Never reuse the sample password or hash as a credential.
\end{description}
Like the SOCKS port, the control port must never be exposed to the network. The control protocol is the subject of \cref{cap:18-stem}.

% ---------------------------------------------------------------------------
\section{Logs}\label{sec:torrc-log}

tor records its activity at several levels of detail: \texttt{err}, \texttt{warn}, \texttt{notice}, \texttt{info}, and \texttt{debug}. The \texttt{notice} level is enough for everyday operation. \texttt{info} and \texttt{debug} are valuable when investigating the protocol, as we did in the laboratories, but produce voluminous output. The \opt{SafeLogging} option, enabled by default, replaces addresses and destinations in logs with \texttt{[scrubbed]}. A forgotten log file must not become a record of the user's activities. The laboratory network deliberately disables this protection (\texttt{SafeLogging 0}) so that we can inspect all the details.

The most useful messages during startup are the \emph{bootstrap} messages. They show the stages of connecting to the network, ending with \texttt{Bootstrapped 100\% (done): Done}. If startup stalls at a particular percentage, the stage suggests where to look. For example, a stall while connecting to the first relay usually indicates censorship or a firewall and suggests using bridges (\cref{cap:16-censura}).

% ---------------------------------------------------------------------------
\section*{Summary}
\begin{itemize}
  \item Use Tor Browser for web browsing; the \texttt{tor} package serves other applications, relays, and onion services.
  \item A \file{torrc} contains one option per line. The manual is the authoritative reference, and \texttt{--verify-config} validates the configuration.
  \item Keep SOCKS and control ports on localhost. Protect the control port with a cookie or password.
  \item Stream isolation prevents unrelated activities from sharing a circuit. Tor Browser achieves this with different SOCKS credentials for each site.
  \item Forcing or excluding relays usually reduces security; \opt{SafeLogging} protects logs.
\end{itemize}

% ---------------------------------------------------------------------------
\section*{Exercises}

\begin{esercizio}[1]\label{es:17-verifica}
In the laboratory container, create a \file{torrc} with a typo in an option name and check it using \texttt{--verify-config}. What message do you get? Then run \texttt{--dump-config short} on the laboratory network client's \file{torrc}.
\end{esercizio}

\begin{esercizio}[2]\label{es:17-isolamento}
A user connects a personal email client under their real identity and a pseudonymous forum through the same \opt{SocksPort}. What could an exit infer? Suggest two configurations that avoid this problem.
\end{esercizio}

\begin{esercizio}[2]\label{es:17-exitnodes}
Why would \texttt{ExitNodes \{it\}}, meaning Italian exits only, reduce the anonymity of an Italian user visiting an Italian website? Reason in terms of the threat model in \cref{cap:04-modello-minaccia}.
\end{esercizio}

% ---------------------------------------------------------------------------
\section*{Further Reading}
\begin{itemize}
  \item The tor manual, \cmd{man tor}, and \file{doc/man/tor.1.txt} in the source tree.
  \item Tor Project support pages on installation and package repositories, \url{https://support.torproject.org/}.
\end{itemize}
